% Structured abstract (ABNT 6028:2021 adaptation) -- max 10 lines.
\begin{abstract}
\textbf{Introduction}: Teaching Nielsen's usability heuristics traditionally relies on reading abstract definitions, which novices often find hard to internalize.
\textbf{Objective}: To support the teaching of Nielsen's heuristics through paired interactive scenarios.
\textbf{Methodology}: Design and evaluation of a web-based educational game where players complete each task in paired violating and compliant interfaces, with narration and gamification; open-ended responses were analysed via qualitative content analysis.
\textbf{Results}: Initial results suggest the paired bad/good contrast is the main mechanism supporting learning over traditional reading; narration served as a focusing aid, gamification as mild engagement scaffolding rather than a learning driver.
\end{abstract}

\keywords{Nielsen's usability heuristics, HCI education, educational game, game-based learning, experience report.}

% Per IHC 2026 rules, a Portuguese "resumo" / "palavras-chave" block is required
% only for papers written in Portuguese. This paper is in English, so the
% resumo and palavraschave blocks are intentionally omitted.
